% c99-contract-test/en.tex — contract test (English, follows de.tex)
%
% Not course content.  This chapter uses every construct of
% docs/authoring/CONTRACT.md at least once and is the fixture for
% scripts/course/check_levels.py and for the HTML converter.  The material
% is deliberately physics-free: arithmetic with integer 2x2 matrices, every
% statement worked out in the text.  It is built only with --chapters c99
% or --with-test-chapter.

\chapter{Contract test: two-by-two matrices}\label{c99}

This chapter is not part of the course. It is a test fixture: it uses every
construct the authoring contract allows, and it computes with integers only.
The course proper begins in Chapter~\ref{c00}; this reference leaves the chapter on
purpose.

\section{Two matrices and their products}\label{c99.s1}

We consider two fixed matrices with integer entries,
\begin{equation}\label{c99.eq.matrizen}\claim{c99.cl.1}{author-stipulated}
  A = \begin{pmatrix} 1 & 2 \\ 0 & 1 \end{pmatrix},
  \qquad
  B = \begin{pmatrix} 1 & 0 \\ 3 & 1 \end{pmatrix}.
\end{equation}
This stipulation is the first of the fixture \srcref{c99-fixture}{STIPULATIONS.md, S1}.
A product of two matrices is formed by the rule “row times column”. For the two
possible orders this gives
\begin{align}
  AB &= \begin{pmatrix} 1\cdot 1 + 2\cdot 3 & 1\cdot 0 + 2\cdot 1 \\
                        0\cdot 1 + 1\cdot 3 & 0\cdot 0 + 1\cdot 1 \end{pmatrix}
      = \begin{pmatrix} 7 & 2 \\ 3 & 1 \end{pmatrix},
  \label{c99.eq.ab}\claim{c99.cl.2}{derived} \\
  BA &= \begin{pmatrix} 1\cdot 1 + 0\cdot 0 & 1\cdot 2 + 0\cdot 1 \\
                        3\cdot 1 + 1\cdot 0 & 3\cdot 2 + 1\cdot 1 \end{pmatrix}
      = \begin{pmatrix} 1 & 2 \\ 3 & 7 \end{pmatrix}.
  \label{c99.eq.ba}\claim{c99.cl.3}{derived}
\end{align}

\begin{proposition}[The order matters]\label{c99.prop.reihenfolge}\claim{c99.cl.4}{derived}
  For the matrices of \eqref{c99.eq.matrizen} we have $AB \neq BA$.
\end{proposition}
\begin{proof}
  The top left entry equals $7$ in \eqref{c99.eq.ab} and $1$ in
  \eqref{c99.eq.ba}.
\end{proof}

\begin{tpcfigure}{c99.fig.scherung}
  \demo{c99.demo.scherung}
  \caption{The matrix $A$ as a map of the plane: the point $(x,y)$ goes to
    $(x+2y,\,y)$. The dashed unit square becomes the hatched parallelogram with
    corners $(0,0)$, $(1,0)$, $(3,1)$ and $(2,1)$; the base stays where it is and
    the top edge moves two units to the right. The German caption of this figure
    is long on purpose, so that the line breaking of long German compounds is
    tested in print and on the web; the English caption is kept at a comparable
    length so that both editions exercise a caption of several lines.}
\end{tpcfigure}

\tpcdef{T1}

With \defref{T1} one reads off from \eqref{c99.eq.ab} and \eqref{c99.eq.ba}:
\begin{equation}\label{c99.eq.spuren}\claim{c99.cl.5}{derived}
  \tr(AB) = 7 + 1 = 8,
  \qquad
  \tr(BA) = 1 + 7 = 8.
\end{equation}
The mean of the four entries of $AB$ is \num[c99.cl.19]{3.25}{}.
The products differ, but their traces agree. This is not an accident of these two
matrices\zoomref{c99.v1}:

\begin{theorem}[Trace of a product]\label{c99.thm.spur}\claim[proof-linked]{c99.cl.6}{derived}
  For all $2\times 2$ matrices $X$ and $Y$ we have $\tr(XY) = \tr(YX)$.
\end{theorem}

\begin{vertiefung}[id=c99.v1,
    titel={Why the traces agree},
    wozu={The core states Theorem~\ref{c99.thm.spur} without proof; here it is
      worked out for arbitrary entries.},
    voraussetzung={Core~\ref{c99.s1} up to equation~\eqref{c99.eq.spuren}; not
      needed for Section~\ref{c99.s2}.},
    zurueck={The trace of a product does not depend on the order of the factors;
      the product itself does.}]
  Let $X$ and $Y$ be matrices with entries $x_{jk}$ and $y_{jk}$. The diagonal
  entries of the two products are
  \begin{align}
    (XY)_{11} &= x_{11}y_{11} + x_{12}y_{21}, &
    (XY)_{22} &= x_{21}y_{12} + x_{22}y_{22},
    \label{c99.eq.diag-xy}\claim{c99.cl.7}{derived} \\
    (YX)_{11} &= y_{11}x_{11} + y_{12}x_{21}, &
    (YX)_{22} &= y_{21}x_{12} + y_{22}x_{22}.
    \label{c99.eq.diag-yx}\claim{c99.cl.8}{derived}
  \end{align}

  \begin{lemma}[Trace of a product, worked out]\label{c99.lem.spur}\claim{c99.cl.9}{derived}
    For all $2\times 2$ matrices $X$ and $Y$ we have
    \begin{equation}\label{c99.eq.spur-allgemein}
      \tr(XY) = x_{11}y_{11} + x_{12}y_{21} + x_{21}y_{12} + x_{22}y_{22} = \tr(YX).
    \end{equation}
  \end{lemma}
  \begin{proof}
    The sum of the two entries in \eqref{c99.eq.diag-xy} is the middle
    expression. The sum of the two entries in \eqref{c99.eq.diag-yx} consists of
    the same four products in a different order.
  \end{proof}

  \tpcdef{T2}

  For the two matrices of the core, \eqref{c99.eq.ab} and \eqref{c99.eq.ba} give
  \begin{equation}\label{c99.eq.kommutator}\claim{c99.cl.10}{derived}
    \comm{A}{B} = \begin{pmatrix} 7-1 & 2-2 \\ 3-3 & 1-7 \end{pmatrix}
                = \begin{pmatrix} 6 & 0 \\ 0 & -6 \end{pmatrix},
    \qquad
    \tr\comm{A}{B} = 6 - 6 = 0.
  \end{equation}

  \begin{tpcfigure}{c99.fig.kommutator}
    \caption{The commutator of \eqref{c99.eq.kommutator} as a map: it stretches
      by the factor $6$ and reflects in the horizontal axis.}
  \end{tpcfigure}

  That the trace respects sums is checked in the foundations\zoomref{c99.f2}.
\end{vertiefung}

The next paragraph is written for readers who skipped the block: it uses only
\eqref{c99.eq.spuren} and Theorem~\ref{c99.thm.spur}.

\begin{tpcfigure}{c99.fig.scherung-b}
  \demo{c99.demo.scherung-b}
  \caption{The matrix $B$ as a map: $(x,y)$ goes to $(x,\,3x+y)$. A second
    figure with a demo on the same page.}
\end{tpcfigure}

\begin{tpctable}{c99.tab.produkte}
  \caption{The entries of the two products of \eqref{c99.eq.ab} and \eqref{c99.eq.ba}.}
  \begin{tabular}{lcccc}
    \toprule
    Product & $(1,1)$ & $(1,2)$ & $(2,1)$ & $(2,2)$ \\
    \midrule
    $AB$ & $7$ & $2$ & $3$ & $1$ \\
    $BA$ & $1$ & $2$ & $3$ & $7$ \\
    \bottomrule
  \end{tabular}
\end{tpctable}

\begin{merksatz}[Key point]
  For matrices the order of the factors matters.
\end{merksatz}

\section{Shears}\label{c99.s2}

\subsection{A family of matrices}\label{c99.s2.1}

For every integer $t$ let
\begin{equation}\label{c99.eq.scherung}\claim{c99.cl.11}{author-stipulated}
  M(t) = \begin{pmatrix} 1 & t \\ 0 & 1 \end{pmatrix}.
\end{equation}
The matrix $A$ of \eqref{c99.eq.matrizen} is $M(2)$.

\begin{theorem}[Shears add]\label{c99.thm.scherung}\claim[proof-included]{c99.cl.12}{derived}
  For all integers $s$ and $t$ we have $M(s)M(t) = M(s+t)$. Moreover, for every
  integer $n \geq 1$ we have $M(t)^n = M(nt)$.
\end{theorem}
\begin{proof}
  The first statement is a multiplication:
  \begin{equation}\label{c99.eq.scherung-produkt}
    M(s)M(t)
    = \begin{pmatrix} 1\cdot 1 + s\cdot 0 & 1\cdot t + s\cdot 1 \\
                      0\cdot 1 + 1\cdot 0 & 0\cdot t + 1\cdot 1 \end{pmatrix}
    = \begin{pmatrix} 1 & s+t \\ 0 & 1 \end{pmatrix}
    = M(s+t).
  \end{equation}
  The second follows by induction on $n$. For $n = 1$ there is nothing to show.
  The induction step is a lemma of its own.
  \begin{lemma}[Induction step]\label{c99.lem.schritt}\claim{c99.cl.13}{derived}
    If $M(t)^n = M(nt)$, then also $M(t)^{n+1} = M((n+1)t)$.
  \end{lemma}
  \begin{proof}
    By assumption $M(t)^{n+1} = M(t)^n M(t) = M(nt)M(t)$, and by
    \eqref{c99.eq.scherung-produkt} with $s = nt$ this is $M(nt + t) = M((n+1)t)$.
  \end{proof}
  With Lemma~\ref{c99.lem.schritt} the statement holds for all $n \geq 1$.
\end{proof}

\begin{corollary}[Shears commute]\label{c99.cor.vertauschen}\claim{c99.cl.18}{derived}
  For all integers $s$ and $t$, $M(s)M(t) = M(t)M(s)$.
\end{corollary}
\begin{proof}
  By Theorem~\ref{c99.thm.scherung} both sides equal $M(s+t)$, because
  $s + t = t + s$.
\end{proof}

\begin{example}[Numbers]\label{c99.ex.zahlen}
  With $s = 2$ and $t = 3$, \eqref{c99.eq.scherung-produkt} gives the equation
  $M(2)M(3) = M(5)$, and with $t = 2$, $n = 3$ we get $A^3 = M(2)^3 = M(6)$.
\end{example}

\begin{remark}
  The set of all $M(t)$ is closed under multiplication. What follows from this
  is in the foundations\zoomref{c99.f1}.
\end{remark}

\subsection{A case distinction}\label{c99.s2.2}

The absolute value of the top right entry of $M(t)$ is
\begin{equation}\label{c99.eq.betrag}\claim{c99.cl.14}{author-stipulated}
  f(t) =
  \begin{cases}
    t  & \text{\tpcifenglish{if}{falls}}\ t \geq 0, \\
    -t & \text{\tpcifenglish{if}{falls}}\ t < 0.
  \end{cases}
\end{equation}

\begin{proposition}[The absolute value is not negative]\label{c99.prop.betrag}\claim{c99.cl.15}{derived}
  For every integer $t$ we have $f(t) \geq 0$ and $f(t)^2 = t^2$.
\end{proposition}
\begin{proof}
  We distinguish the two cases of \eqref{c99.eq.betrag}.
  \begin{description}
    \item[Case 1: $t \geq 0$.] Then $f(t) = t \geq 0$ and $f(t)^2 = t^2$.
    \item[Case 2: $t < 0$.] Then $f(t) = -t > 0$ and
      $f(t)^2 = (-t)^2 = t^2$.
  \end{description}
\end{proof}

\begin{intuition}
  A shear moves every horizontal line along itself; only the amount of the shift
  grows with the height. The boxes for element types are tinted, the depth
  blocks are not.
\end{intuition}

\begin{exercise}[The inverse]\label{c99.exc.umkehrung}
  Compute $M(t)M(-t)$.
  \begin{enumerate}
    \item Use \eqref{c99.eq.scherung-produkt}.
    \item Compare with the direct multiplication.
  \end{enumerate}
\end{exercise}
\begin{solution}
  With $s = t$ and $-t$ in the place of $t$, \eqref{c99.eq.scherung-produkt}
  gives the matrix $M(t - t) = M(0)$, and $M(0)$ has the entries $1, 0, 0, 1$.
  \begin{itemize}
    \item Directly: the top right entry is $1\cdot(-t) + t\cdot 1 = 0$.
    \item The remaining entries are as in \eqref{c99.eq.scherung-produkt}.
  \end{itemize}
\end{solution}

\begin{fundament}[id=c99.f1,
    titel={The shears form a group},
    wozu={The core uses only the product of two shears; here we collect what
      follows from it for the whole family.},
    voraussetzung={Core~\ref{c99.s2.1}, Theorem~\ref{c99.thm.scherung}.},
    zurueck={Every shear $M(t)$ is undone by $M(-t)$.}]
  \begin{proposition}[Identity and inverse]\label{c99.prop.gruppe}\claim{c99.cl.16}{derived}
    For all integers $t$ we have $M(0)M(t) = M(t)$ and $M(t)M(-t) = M(0)$,
    and likewise with the factors exchanged.
  \end{proposition}
  \begin{proof}
    Both are Theorem~\ref{c99.thm.scherung} with a special choice of $s$:
    \begin{equation}\label{c99.eq.gruppe}
      M(0)M(t) = M(0 + t) = M(t),
      \qquad
      M(t)M(-t) = M(t - t) = M(0).
    \end{equation}
    The equations with the factors exchanged follow in the same way, because
    $s + t = t + s$.
  \end{proof}
\end{fundament}

\begin{fundament}[id=c99.f2,
    titel={The trace respects sums},
    wozu={The deeper block uses that the trace of a difference is the difference
      of the traces; here is the computation.},
    voraussetzung={\defref{T1}.},
    zurueck={From $\tr(AB) = \tr(BA)$ it therefore follows that $\tr(AB - BA) = 0$.}]
  For matrices $X$ and $Y$ with entries $x_{jk}$ and $y_{jk}$ and an integer $c$,
  \begin{equation}\label{c99.eq.spur-linear}\claim{c99.cl.17}{derived}
    \tr(X + cY) = (x_{11} + c\,y_{11}) + (x_{22} + c\,y_{22})
                = \tr X + c\,\tr Y.
  \end{equation}
  With $c = -1$ this is the statement about differences.
\end{fundament}
